% Shared content style for NN and NNDS. Written by Claude (2026-10-14).
%
% Everything that shared slides rely on lives here, loaded by both courses:
% colors, notation helpers and the card system. Course themes own only the
% presentation around the content (header, footer, title page, fonts, the
% \myalert color and where \SlideSource text goes). Change a definition here,
% never by redefining it in a course file, so the two courses cannot drift.

\usepackage{mathtools,amssymb,tcolorbox}
\tcbsetforeverylayer{autoparskip}

% Semantic colors
\definecolor{darkred}{rgb}{0.55,0,0}
\definecolor{darkpastelblue}{rgb}{0.47,0.62,0.8}
\definecolor{darkcerulean}{rgb}{0.03,0.27,0.49}
\definecolor{darkgreen}{rgb}{0,0.55,0}

% Figure palette (draw.io defaults)
\definecolor{titlecolor}{RGB}{245, 242, 238}
\definecolor{drawgray}{HTML}{666666}
\definecolor{drawgray_l}{HTML}{F5F5F5}
\definecolor{drawblue}{HTML}{6C8EBF}
\definecolor{drawblue_l}{HTML}{DAE8FC}
\definecolor{drawgreen}{HTML}{82B366}
\definecolor{drawgreen_l}{HTML}{D5E8D4}
\definecolor{draworange}{HTML}{D79B00}
\definecolor{draworange_l}{HTML}{FFE6CC}
\definecolor{drawyellow}{HTML}{D6B656}
\definecolor{drawyellow_l}{HTML}{FFF2CC}
\definecolor{drawred}{HTML}{B85450}
\definecolor{drawred_l}{HTML}{EA6B66}
\definecolor{drawviolet}{HTML}{9673A6}
\definecolor{drawviolet_l}{HTML}{E1D5E7}

% Notation
\newcommand{\vc}[1]{\boldsymbol{\mathbf{#1}}}
\DeclareMathOperator*{\argmax}{arg\,max}
\DeclareMathOperator*{\argmin}{arg\,min}

% Cards: blue for definitions and computable quantities, green for goals and
% stable formulations, orange for warnings and trade-offs, red for
% training-specific or problematic elements.
\tcbset{
	nndsblue/.style={colback=drawblue_l!35!white,colframe=drawblue,boxrule=0.8pt,arc=1mm},
	nndsgreen/.style={colback=drawgreen_l!35!white,colframe=drawgreen,boxrule=0.8pt,arc=1mm},
	nndsorange/.style={colback=draworange_l!35!white,colframe=draworange,boxrule=0.8pt,arc=1mm},
	nndsred/.style={colback=drawred_l!12!white,colframe=drawred,boxrule=0.8pt,arc=1mm}
}

% A card that opens with \[ would start with an empty paragraph line plus the
% display skip, leaving more room above than below (by an amount that differs
% between engines). Set that leading equation as a centered display-style line
% instead, so the space above it is the card padding, as below.
\makeatletter
\newcommand{\nndscardlead}{\futurelet\nndscard@next\nndscard@check}
% tcolorbox follows this hook with \ignorespaces; look past it.
\def\nndscard@check{\ifx\nndscard@next\ignorespaces
	\expandafter\nndscard@skip\else\expandafter\nndscard@test\fi}
\def\nndscard@skip\ignorespaces{\nndscard@test}
\newcommand{\nndscard@test}{\@ifnextchar\[{\nndscard@display}{}}
\def\nndscard@display\[#1\]{{\centering$\displaystyle#1$\par}\@ifnextchar\end{}{\vspace{\belowdisplayskip}}}
\makeatother

% Card text is ragged-right: narrow justified columns open wide word gaps,
% especially in NNDS, which never hyphenates. Fixed display spacing: when
% parallel cards share a height, stretchable skips would otherwise spread
% the content unevenly.
\tcbset{
	nndscard/.style={before upper={\raggedright
		\setlength{\abovedisplayskip}{6pt}\setlength{\belowdisplayskip}{6pt}%
		\setlength{\abovedisplayshortskip}{6pt}\setlength{\belowdisplayshortskip}{6pt}%
		\nndscardlead}},
	nndsblue/.append style={nndscard},
	nndsgreen/.append style={nndscard},
	nndsorange/.append style={nndscard},
	nndsred/.append style={nndscard}
}

% Parallel cards share a group name and get equal heights; the second argument
% (a fixed height used by older NNDS slides) is ignored.
\tcbset{courseheight/.style 2 args={equal height group=#1}}

\usepackage{tikz}

% Matrix shapes written under a symbol, e.g. \shape{\vc{Q}}{(T,D_k)}.
\usepackage{stackengine}
\stackMath
\def\sss{\scriptstyle \color{gray}}
\setstackgap{L}{12pt}
\def\stacktype{L}
\newcommand{\shape}[2]{\stackunder{#1}{\sss #2}}

% Transformer slides: colors and diagram styles (moved here from the NNDS
% transformers-settings.tex on 2026-11-16, when NN adopted the lecture).

% Query, key and value colors, reused in equations and figures.
\colorlet{qcolor}{drawviolet!80!black}
\colorlet{kcolor}{draworange!85!black}
\colorlet{vcolor}{darkcerulean}
\newcommand{\qc}[1]{{\color{qcolor}#1}}
\newcommand{\kc}[1]{{\color{kcolor}#1}}
\newcommand{\vcl}[1]{{\color{vcolor}#1}}

% Block-diagram styles, as in shared/slides/mlps.tex and the deeper-models deck.
\tikzset{
    nnblock/.style={
            draw=drawblue, fill=drawblue_l!25!white, rounded corners=2pt,
            minimum width=1.25cm, minimum height=0.65cm,
            inner sep=5pt, font=\small
        },
    nnother/.style={nnblock, draw=drawgreen, fill=drawgreen_l!25!white},
    nnmerge/.style={
            circle, draw=drawgray, fill=white,
            minimum size=0.55cm, inner sep=1pt, font=\small
        },
    nnedge/.style={->, thick, draw=drawgray}
}

% Full-size diagram styles for the transformer figures: arrows as in the
% attention figures (Stealth tips, gray), boxes at normal size.
\tikzset{
    tfarr/.style={-{Stealth[length=5pt]}, drawgray, line width=0.6pt},
    tfblk/.style={draw=drawblue, fill=drawblue_l!60, rounded corners=2pt,
            minimum height=0.75cm, minimum width=1.9cm, inner sep=4pt, font=\small},
    tfnorm/.style={tfblk, draw=drawgreen, fill=drawgreen_l!70},
    tfsum/.style={circle, draw=drawgray, fill=white, minimum size=0.45cm, inner sep=0pt, font=\small},
    tflab/.style={font=\scriptsize, text=black!60},
}
